\documentclass[10pt,a4paper]{amsart}
\usepackage{amsmath,amssymb,amsthm,graphicx}
% Retained preamble for the Stacks Project chapter snapshot in this directory.
% The chapter's own preamble.tex is a copy of the shared Stacks Project preamble; the block
% below is sliced MECHANICALLY, line for line, from it. Nothing was added, removed or
% redefined; the wrapper class, the xr-hyper/hyperref wiring and the externaldocument list
% are omitted because the assembler supplies the document class.
%
% lines taken from preamble.tex: 3, 5, 9, 11, 12, 156, 157, 158, ... (51 lines total)
%
% package coverage restored from the chapter's own preamble (2026-10-13): xy, xy-2cell

\usepackage[all]{xy}
\xyoption{2cell}
\UseAllTwocells

\newenvironment{reference}{\comment}{\endcomment}
\newenvironment{slogan}{\comment}{\endcomment}
\newenvironment{history}{\comment}{\endcomment}
\theoremstyle{plain}
\newtheorem{theorem}[subsection]{Theorem}
\newtheorem{proposition}[subsection]{Proposition}
\newtheorem{lemma}[subsection]{Lemma}
\theoremstyle{definition}
\newtheorem{definition}[subsection]{Definition}
\newtheorem{example}[subsection]{Example}
\newtheorem{exercise}[subsection]{Exercise}
\newtheorem{situation}[subsection]{Situation}
\theoremstyle{remark}
\newtheorem{remark}[subsection]{Remark}
\newtheorem{remarks}[subsection]{Remarks}
\numberwithin{equation}{subsection}
\def\lim{\mathop{\mathrm{lim}}\nolimits}
\def\colim{\mathop{\mathrm{colim}}\nolimits}
\def\Spec{\mathop{\mathrm{Spec}}}
\def\Hom{\mathop{\mathrm{Hom}}\nolimits}
\def\Ext{\mathop{\mathrm{Ext}}\nolimits}
\def\SheafHom{\mathop{\mathcal{H}\!\mathit{om}}\nolimits}
\def\SheafExt{\mathop{\mathcal{E}\!\mathit{xt}}\nolimits}
\def\Sch{\mathit{Sch}}
\def\Mor{\mathop{\mathrm{Mor}}\nolimits}
\def\Ob{\mathop{\mathrm{Ob}}\nolimits}
\def\Sh{\mathop{\mathit{Sh}}\nolimits}
\def\NL{\mathop{N\!L}\nolimits}
\def\CH{\mathop{\mathrm{CH}}\nolimits}
\def\proetale{{pro\text{-}\acute{e}tale}}
\def\etale{{\acute{e}tale}}
\def\QCoh{\mathit{QCoh}}
\def\Ker{\mathop{\mathrm{Ker}}}
\def\Im{\mathop{\mathrm{Im}}}
\def\Coker{\mathop{\mathrm{Coker}}}
\def\Coim{\mathop{\mathrm{Coim}}}
\DeclareMathSymbol{\boxtimes}{\mathbin}{AMSa}{"02}
\def\QCohstack{\mathcal{QC}\!\mathit{oh}}
\def\Cohstack{\mathcal{C}\!\mathit{oh}}
\def\Spacesstack{\mathcal{S}\!\mathit{paces}}
\def\Quotfunctor{\mathrm{Quot}}
\def\Hilbfunctor{\mathrm{Hilb}}
\def\Curvesstack{\mathcal{C}\!\mathit{urves}}
\def\Polarizedstack{\mathcal{P}\!\mathit{olarized}}
\def\Complexesstack{\mathcal{C}\!\mathit{omplexes}}
\def\Pic{\mathop{\mathrm{Pic}}\nolimits}
\def\Picardstack{\mathcal{P}\!\mathit{ic}}
\def\Picardfunctor{\mathrm{Pic}}
\def\Deformationcategory{\mathcal{D}\!\mathit{ef}}
\usepackage{cleveref}
\crefname{originallemma}{original lemma}{original lemmas}
\Crefname{originallemma}{Original lemma}{Original lemmas}
\makeatletter
\newlabel{divisors-lemma-base-change-and-fitting-ideal-omega}{{0C3I}{source}{}{sourcepointer-1}{}}
\newlabel{divisors-lemma-base-change-and-fitting-ideal-omega@cref}{{[originallemma][0][]0C3I}{[1][0][]source}{}{}{}}
\makeatother
\crefname{originallemma}{original lemma}{original lemmas}
\Crefname{originallemma}{Original lemma}{Original lemmas}
\makeatletter
\newlabel{divisors-lemma-d-fitting-ideal-omega-smooth}{{0C3K}{source}{}{sourcepointer-2}{}}
\newlabel{divisors-lemma-d-fitting-ideal-omega-smooth@cref}{{[originallemma][0][]0C3K}{[1][0][]source}{}{}{}}
\makeatother
\crefname{originallemma}{original lemma}{original lemmas}
\Crefname{originallemma}{Original lemma}{Original lemmas}
\makeatletter
\newlabel{more-morphisms-lemma-flat-finite-presentation-CM-open}{{045U}{source}{}{sourcepointer-3}{}}
\newlabel{more-morphisms-lemma-flat-finite-presentation-CM-open@cref}{{[originallemma][0][]045U}{[1][0][]source}{}{}{}}
\makeatother
\crefname{originallemma}{original lemma}{original lemmas}
\Crefname{originallemma}{Original lemma}{Original lemmas}
\makeatletter
\newlabel{morphisms-lemma-flat-finite-presentation-CM-fibres-relative-dimension}{{02NM}{source}{}{sourcepointer-4}{}}
\newlabel{morphisms-lemma-flat-finite-presentation-CM-fibres-relative-dimension@cref}{{[originallemma][0][]02NM}{[1][0][]source}{}{}{}}
\makeatother
\title{Stacks, Examples of Deformations: lemma-criterion-smoothing}
\author{The Stacks Project Authors}
\date{stacks-project:examples-defos.tex:local-snapshot-20261012; GNU FDL 1.2 or later}
\begin{document}
\maketitle
\noindent\textit{Editorial note.} Statement span is lines 3116-3124 of sources/stacks-examples-defos/pinned-original-source.tex; proof span is lines 3126-3148, the complete author proof (21 body lines).\par\medskip
\noindent\textit{Editorial note.} Navigation stubs: 4 cited-but-not-included labels are supplied as non-load-bearing pointers with their pinned tags where the snapshot tag index has them.\par\medskip
\noindent\textit{Editorial note.} extra\_packages is empty on purpose: the retained preamble already loads xy with options and a second declaration clashes in the assembler wrapper.\par\medskip
\noindent\textit{Editorial note.} Text is retained verbatim, including the author's own typos and placeholders.\par\medskip
\noindent\textit{Editorial note.} Queue index 23 of queue\_examples-defos.json (pb=1121, stmt 3116-3124, proof 3126-3148).\par\medskip
\noindent\textit{Allowed standard prerequisites (editorial).}\begin{itemize}
\item Explicit prerequisite boundary. The cited inputs are prerequisites; the argument in the retained span is not replaced by any cited result.
\end{itemize}\medskip
\ifdefined\hypertarget\hypertarget{sourcepointer-1}{}\fi
\noindent\textit{Source reference.} Original lemma 0C3I: sources/stacks-examples-defos/pinned-original-source.tex, source lines 3148--3148. Cited by the retained statement or proof but not part of the collected spans; supplied only as a navigation pointer. The argument does not depend on its content.\par\medskip
\ifdefined\hypertarget\hypertarget{sourcepointer-2}{}\fi
\noindent\textit{Source reference.} Original lemma 0C3K: sources/stacks-examples-defos/pinned-original-source.tex, source lines 3148--3148. Cited by the retained statement or proof but not part of the collected spans; supplied only as a navigation pointer. The argument does not depend on its content.\par\medskip
\ifdefined\hypertarget\hypertarget{sourcepointer-3}{}\fi
\noindent\textit{Source reference.} Original lemma 045U: sources/stacks-examples-defos/pinned-original-source.tex, source lines 3148--3148. Cited by the retained statement or proof but not part of the collected spans; supplied only as a navigation pointer. The argument does not depend on its content.\par\medskip
\ifdefined\hypertarget\hypertarget{sourcepointer-4}{}\fi
\noindent\textit{Source reference.} Original lemma 02NM: sources/stacks-examples-defos/pinned-original-source.tex, source lines 3148--3148. Cited by the retained statement or proof but not part of the collected spans; supplied only as a navigation pointer. The argument does not depend on its content.\par\medskip
\begin{lemma}
\label{lemma-criterion-smoothing}
Let $k$ be a field. Set $S = \Spec(k[[t]])$ and
$S_n = \Spec(k[t]/(t^n))$. Let $Y \to S$ be a proper, flat morphism
of schemes whose special fibre $X$ is Cohen-Macaulay and
equidimensional of dimension $d$. Denote $X_n = Y \times_S S_n$.
If for some $n \geq 1$ the $d$th Fitting ideal of $\Omega_{X_n/S_n}$
contains $t^{n - 1}$, then the generic fibre of $Y \to S$ is smooth.
\end{lemma}

\begin{proof}
By More on Morphisms, Lemma
\ref{more-morphisms-lemma-flat-finite-presentation-CM-open}
we see that $Y \to S$ is a Cohen-Macaulay morphism.
By Morphisms, Lemma
\ref{morphisms-lemma-flat-finite-presentation-CM-fibres-relative-dimension}
we see that $Y \to S$ has relative dimension $d$.
By Divisors, Lemma \ref{divisors-lemma-d-fitting-ideal-omega-smooth}
the $d$th Fitting ideal $\mathcal{I} \subset \mathcal{O}_Y$
of $\Omega_{Y/S}$ cuts out the singular locus of the morphism $Y \to S$.
In other words, $V(\mathcal{I}) \subset Y$ is the closed subset
of points where $Y \to S$ is not smooth.
By Divisors, Lemma \ref{divisors-lemma-base-change-and-fitting-ideal-omega}
formation of this Fitting ideal commutes with base change.
By assumption we see that $t^{n - 1}$ is a section of
$\mathcal{I} + t^n\mathcal{O}_Y$. Thus for every
$x \in X = V(t) \subset Y$ we conclude that
$t^{n - 1} \in \mathcal{I}_x$ where $\mathcal{I}_x$ is the stalk at $x$.
This implies that $V(\mathcal{I}) \subset V(t)$ in an
open neighbourhood of $X$ in $Y$. Since $Y \to S$
is proper, this implies $V(\mathcal{I}) \subset V(t)$
as desired.
\end{proof}

\end{document}
